% =====================================================================
%  Dokumentation: Projekt "Goldene Mitte"
%  Übersetzen mit:  pdflatex GoldeneMitte.tex   (zweimal, für das Inhaltsverzeichnis)
% =====================================================================
\documentclass[11pt,a4paper,parskip=half]{scrartcl}

\usepackage[T1]{fontenc}
\usepackage[ngerman]{babel}
\usepackage{lmodern}
\usepackage{microtype}
\usepackage[autostyle]{csquotes}
\usepackage[margin=2.5cm]{geometry}
\usepackage[table]{xcolor}
\usepackage{booktabs}
\usepackage{tabularx}
\usepackage{float}
\usepackage{amsmath}
\usepackage{listings}
\usepackage{algorithm}
\usepackage{algpseudocode}
\usepackage{tikz}
\usetikzlibrary{shapes.multipart, positioning, arrows.meta, calc, fit, backgrounds}
\usepackage[section,above]{placeins}
\usepackage{hyperref}
\setlength{\emergencystretch}{3em}

\hypersetup{colorlinks=true, linkcolor=blue!50!black, urlcolor=blue!50!black,
            pdftitle={Goldene Mitte -- Projektdokumentation}}

% --------------------------------------------------------------------
%  Allgemeine Makros
% --------------------------------------------------------------------
\definecolor{goldcol}{HTML}{D4AF37}
\newcommand{\gold}{\cellcolor{goldcol}\textbf{G}}
\newcommand{\code}[1]{\texttt{#1}}
\newcommand{\file}[1]{\texttt{#1}}

\lstset{
  basicstyle=\ttfamily\small,
  columns=fullflexible,
  keepspaces=true,
  frame=single,
  framesep=4pt,
  backgroundcolor=\color{gray!6},
  rulecolor=\color{gray!40},
  showstringspaces=false,
  literate={ä}{{\"a}}1 {ö}{{\"o}}1 {ü}{{\"u}}1 {Ä}{{\"A}}1 {Ö}{{\"O}}1 {Ü}{{\"U}}1 {ß}{{\ss}}1
}
\lstdefinestyle{java}{language=Java, keywordstyle=\color{blue!60!black}\bfseries,
  commentstyle=\color{green!40!black}\itshape, stringstyle=\color{red!50!black}}

\algrenewcommand\algorithmicrequire{\textbf{Eingabe:}}
\algrenewcommand\algorithmicensure{\textbf{Ausgabe:}}
\floatname{algorithm}{Algorithmus}

% --------------------------------------------------------------------
%  UML-Klassenboxen (TikZ)
%   \umlclass[Positionierung]{Knotenname}{Titel}{Attribute}{Methoden}
% --------------------------------------------------------------------
\tikzset{
  umlclass/.style={draw, rectangle split, rectangle split parts=3,
    rectangle split part fill={blue!8, white, white},
    text width=#1, font=\ttfamily\scriptsize, align=left, inner sep=3pt},
  umlclass/.default=6.2cm,
  dep/.style={-{Stealth[length=2.2mm]}, dashed, thin},
  assoc/.style={-{Stealth[length=2.2mm]}, thin},
  comp/.style={{Diamond[length=3mm,width=2mm]}-{Stealth[length=2.2mm]}, thin},
  aggr/.style={{Diamond[open,length=3mm,width=2mm]}-{Stealth[length=2.2mm]}, thin},
  lbl/.style={font=\ttfamily\scriptsize, fill=white, inner sep=1pt},
}
\newcommand{\umlclass}[5][]{%
  \node[umlclass, #1] (#2)
    {\centering\textbf{#3}\par
     \nodepart{second}#4
     \nodepart{third}#5};}
\newcommand{\stereo}[1]{{\normalfont\guillemotleft}#1{\normalfont\guillemotright}}
\newcommand{\st}[1]{\underline{#1}} % statische Elemente unterstrichen

% =====================================================================
\title{Goldene Mitte}
\subtitle{Projektdokumentation -- Lernfeld 13}
\author{\textit{Name eintragen}}
\date{\today}

\begin{document}
\maketitle
\tableofcontents
\newpage

% =====================================================================
\section{Einführung}
% =====================================================================

\subsection{Aufgabenstellung}

Dirk hat zum Geburtstag ein Rätsel bekommen: eine würfelförmige Kiste mit einer
inneren Kantenlänge von 5\,cm und eine Menge von Würfeln und Quadern. Darunter
befindet sich ein \emph{goldener Würfel} der Kantenlänge 1\,cm. Die Teile sollen
so in die Kiste gelegt werden, dass diese \emph{vollständig} gefüllt ist und der
goldene Würfel \emph{genau in der Mitte} liegt.

Neben dem goldenen Würfel stehen zur Verfügung:
\begin{itemize}
  \item vier Würfel mit Kantenlänge 1\,cm,
  \item sechs Quader $1 \times 2 \times 4$\,cm,
  \item sechs Quader $2 \times 2 \times 3$\,cm.
\end{itemize}

Gesucht ist ein Programm, das die Kiste nach diesen Regeln füllt oder meldet, dass es
keine Lösung gibt. Das Programm soll auch für andere Kisten und andere Mengen von
Teilen funktionieren. Dabei gilt immer:
\begin{itemize}
  \item alle Kantenlängen sind ganzzahlig,
  \item ein goldener $1\times1\times1$-Würfel ist immer dabei und liegt in der Mitte.
\end{itemize}

\subsection{Grundlegende Beobachtungen}

\begin{itemize}
  \item Weil der goldene Würfel \emph{genau} in der Mitte liegen soll, muss die
        Kantenlänge $n$ der Kiste \textbf{ungerade} sein. Nur dann gibt es eine
        einzelne Mittelzelle $\left(\lfloor n/2\rfloor, \lfloor n/2\rfloor,
        \lfloor n/2\rfloor\right)$.
  \item Da die Kantenlängen ganzzahlig sind, kann man die Kiste als
        dreidimensionales Gitter aus $n^3$ Einheitszellen auffassen. Jedes Teil
        belegt einen achsenparallelen Block dieser Zellen.
  \item Eine Lösung kann es nur geben, wenn das Gesamtvolumen der Teile genau
        $n^3 - 1$ beträgt, also allen Zellen außer der Mittelzelle entspricht.
  \item Jedes Teil darf gedreht werden. Ein Quader $a\times b\times c$ hat höchstens
        $3! = 6$ verschiedene achsenparallele Lagen.
\end{itemize}

\subsection{Projektüberblick}

Das Projekt ist in Java~25 mit Maven umgesetzt. Es besteht aus zwei Teilen:
\begin{description}
  \item[\code{GoldeneMitte.core}] Einlesen der Rätseldateien, Datenmodell
        (Kiste, Baustein, Position) und der Lösungsalgorithmus
        (Backtracking).
  \item[\code{GoldeneMitte.visualizer}] Eine JavaFX-Oberfläche nach dem
        Model-View-Controller-Muster, die die gefundene Lösung dreidimensional
        darstellt. Die Ansicht lässt sich drehen und zoomen, und einzelne Bausteine
        können ausgewählt werden.
\end{description}

% =====================================================================
\section{Benutzeranleitung}
% =====================================================================

\subsection{Voraussetzungen}

\begin{itemize}
  \item JDK~25 (die \file{pom.xml} kompiliert mit \code{maven.compiler.release = 25}),
  \item Apache Maven~3.9 oder neuer, alternativ IntelliJ IDEA mit Maven-Unterstützung,
  \item JavaFX~21 und JUnit~5. Beide werden automatisch von Maven geladen.
\end{itemize}

\textbf{Hinweis:} Maven muss mit einem JDK~25 laufen. Ist ein älteres JDK
eingestellt, bricht der Build ab. Dann muss \code{JAVA\_HOME} vorher auf das JDK~25
gesetzt werden.

\subsection{Format der Eingabedateien}

Die Rätsel liegen als Textdateien im Ordner \file{src/main/resources/input/}. Jede
Datei enthält
\begin{enumerate}
  \item in der ersten Zeile die inneren Kantenlängen der Kiste: \code{x y z},
  \item in der zweiten Zeile die Anzahl $k$ der Teile,
  \item in den folgenden $k$ Zeilen jeweils die Kantenlängen \code{a b c} eines Teils.
\end{enumerate}
Der goldene Würfel wird \emph{nicht} aufgeführt, er ist immer dabei. Die Datei für
Dirks Rätsel (\file{raetsel5.txt}) beginnt zum Beispiel so:

\begin{lstlisting}
5 5 5
16
1 1 1
2 2 3
2 2 3
...
\end{lstlisting}

Mitgeliefert werden die Beispiele \file{raetsel1.txt} bis \file{raetsel5.txt} sowie
zwei absichtlich fehlerhafte Rätsel, \file{raetsel\_nicht\_lösbar.txt} und
\file{raetsel\_unvollständig.txt} (siehe Abschnitt~\ref{sec:ergebnisse}).

\subsection{Programm starten}

Es gibt zwei Einstiegspunkte:

\begin{description}
  \item[Konsolenversion -- \code{GoldeneMitte.Main}]
    Liest ein Rätsel ein, löst es und gibt das Ergebnis (\code{true}/\code{false})
    sowie die Kiste Ebene für Ebene auf der Konsole aus. Start aus dem
    Projektverzeichnis:
\begin{lstlisting}
mvn compile
java -cp target/classes GoldeneMitte.Main
\end{lstlisting}

  \item[Grafische Version -- \code{GoldeneMitte.Launcher}]
    Startet die JavaFX-Anwendung \code{Visualizer}. Diese löst das Rätsel, schreibt
    die Lösung zusätzlich in die Datei \file{src/main/resources/\allowbreak output/\allowbreak raetsel4.txt}
    und zeigt sie in 3D an. In IntelliJ wird dazu die Klasse \code{Launcher} über den
    grünen Pfeil gestartet. Die eigene Klasse \code{Launcher} ist nötig, weil der
    JavaFX-Start aus einer Klasse, die selbst von \code{Application} erbt, ohne
    Modulpfad fehlschlägt.
\end{description}

Beide Programme müssen mit dem \emph{Projektverzeichnis} als Arbeitsverzeichnis
gestartet werden, weil die Dateipfade relativ angegeben sind.

\paragraph{Anderes Rätsel auswählen.}
Der Pfad der Eingabedatei steht zurzeit fest im Quelltext, jeweils in der ersten Zeile
von \code{Main.main} bzw. \code{Visualizer.start}:
\begin{lstlisting}[style=java]
PuzzleInput input = PuzzleInput.fromFile("src/main/resources/input/raetsel4.txt");
\end{lstlisting}
Für ein anderes Rätsel wird der Dateiname hier geändert.

\subsection{Ausgabe lesen}

Die Konsolenausgabe zeigt die Kiste in Schichten (\emph{Layer}) entlang der
$z$-Achse. In jeder Schicht steht eine Zeile für einen $y$-Wert und eine Spalte für
einen $x$-Wert. Jede Zahl ist die ID des Teils, das diese Zelle belegt. \code{G}
markiert den goldenen Würfel, ein Punkt eine leere Zelle.

\textbf{Wichtig:} Die IDs beginnen bei \textbf{0}. ID~$i$ steht für das Teil in
Zeile $i+3$ der Eingabedatei, also das $(i+1)$-te aufgeführte Teil.

\subsection{Bedienung der 3D-Ansicht}

\begin{center}
\begin{tabularx}{\textwidth}{@{}lX@{}}
\toprule
\textbf{Aktion} & \textbf{Wirkung} \\
\midrule
Linke Maustaste ziehen & Kiste drehen \\
Mausrad & Zoomen \\
Klick auf einen Baustein (3D oder Seitenleiste) & Nur diesen Baustein farbig anzeigen. Die übrigen erscheinen als graue Umrisse. \\
Erneuter Klick auf denselben Baustein & Auswahl aufheben \\
Klick ins Leere & Auswahl aufheben \\
Schaltfläche \enquote{Ansicht zurücksetzen} & Kamera in die Ausgangslage bringen \\
\bottomrule
\end{tabularx}
\end{center}

Die Seitenleiste listet alle Bausteine mit Farbe, ID und Anzahl der belegten Zellen
(\enquote{Voxel}) auf. Wurde keine Lösung gefunden, trägt das Fenster den Titel
\enquote{Demo (nicht gelöst)}.

% =====================================================================
\section{Algorithmus}
% =====================================================================

Das Programm löst das Rätsel mit \textbf{Backtracking}, also einer systematischen
Tiefensuche über alle Möglichkeiten, die Teile zu legen. Mehrere Maßnahmen verkleinern
den Suchraum so stark, dass auch die $5\times5\times5$-Rätsel in wenigen Sekunden
gelöst werden.

\subsection{Datenmodell}

\begin{itemize}
  \item Die Kiste (\code{Box}) ist ein Feld \code{Brick[n][n][n]}. Eine Zelle ist
        \code{null}, wenn sie frei ist, und enthält sonst eine Referenz auf den
        Baustein, der sie belegt. Der Konstruktor legt sofort den goldenen Würfel in
        die Mittelzelle und lehnt gerade oder nicht positive Kantenlängen mit einer
        \code{IllegalArgumentException} ab.
  \item Ein Baustein (\code{Brick}) speichert seine ID, die Kantenlängen
        $x, y, z$ und die Liste seiner \emph{verschiedenen} Orientierungen.
\end{itemize}

\subsection{Schritt 1: Volumenprüfung}

Vor der eigentlichen Suche wird geprüft, ob
\[
  \sum_{i} a_i \cdot b_i \cdot c_i \;=\; n^3 - 1 .
\]
Ist das nicht der Fall, ist das Rätsel sicher unlösbar und das Programm gibt sofort
\code{false} zurück (Methode \code{unsolvableFromStart}). So werden zum Beispiel die
Dateien \file{raetsel\_nicht\_lösbar.txt} (Volumen~28 statt~26) und
\file{raetsel\_unvollständig.txt} (Volumen~22 statt~26) ohne Suche abgelehnt.

\subsection{Schritt 2: Orientierungen ohne Duplikate}

Ein Quader kann auf $3! = 6$ Arten achsenparallel gedreht werden: Jede Permutation
der Kantenlängen auf die Achsen $x, y, z$ ergibt eine Lage. Diese sechs Permutationen
stehen in \code{Puzzle.orientations}. Bei Teilen mit gleichen Kantenlängen ergeben
einige Permutationen dieselbe Lage. Ein Würfel hat zum Beispiel nur eine einzige.
Beim Erzeugen eines \code{Brick} werden die resultierenden Maße deshalb als
\code{Triple} in einem \code{HashSet} gesammelt, und es bleiben nur die
verschiedenen Lagen übrig:

\begin{center}
\begin{tabular}{@{}lc@{\hspace{2.5em}}lc@{}}
\toprule
\textbf{Teil} & \textbf{Lagen} & \textbf{Teil} & \textbf{Lagen} \\
\midrule
$1\times1\times1$, $2\times2\times2$ & 1 & $1\times2\times3$, $1\times2\times4$ & 6 \\
$1\times1\times2$, $1\times1\times5$, $1\times2\times2$ & 3 & $2\times2\times3$, $1\times3\times3$ & 3 \\
\bottomrule
\end{tabular}
\end{center}

\subsection{Schritt 3: Gleiche Teile zusammenfassen}

Viele Rätsel enthalten mehrere gleiche Teile, Dirks Rätsel etwa sechsmal
$2\times2\times3$. Würde man jedes Teil einzeln betrachten, probierte die Suche jede
Vertauschung gleicher Teile erneut aus, obwohl das zur selben Belegung führt.

Deshalb gelten zwei \code{Brick}-Objekte als \emph{gleich}, wenn ihre
\emph{sortierten} Kantenlängen übereinstimmen (\code{equals}/\code{hashCode}
vergleichen \code{sortedDimensions()}). Ein $1\times2\times4$- und ein
$4\times1\times2$-Teil sind also gleich. Die Methode \code{distinctBricks} gruppiert
die Teile in eine \code{Map<Brick, Set<Integer>>} von \emph{Teiltyp} auf die IDs der
noch verfügbaren Exemplare. Die Suche wählt dann nur noch einen \emph{Typ} aus und
nimmt irgendein freies Exemplar davon.

Bei \file{raetsel4.txt} schrumpft das von 20 einzelnen Teilen auf 5 Typen. Die
Anzahl der eingesparten gleichwertigen Vertauschungen beträgt
$4!\cdot 9!\cdot 5! \approx 10^9$.

\subsection{Schritt 4: Füllen der ersten freien Zelle}

Die Zellen werden in einer festen Reihenfolge durchlaufen: zuerst $x$, dann $y$, dann
$z$ (Methode \code{nextFreeCell}). Die Suche betrachtet immer die \emph{erste freie
Zelle} $c$ in dieser Reihenfolge und legt ein Teil so, dass seine \emph{kleinste
Ecke} genau in $c$ liegt.

\paragraph{Warum reicht das?}
Irgendein Teil muss $c$ in einer Lösung belegen. Seine kleinste Ecke $m$ liegt in
jeder Koordinate höchstens bei $c$, kommt in der Durchlaufreihenfolge also nicht nach
$c$. Wäre $m \neq c$, so läge $m$ \emph{vor} $c$. Dann wäre $m$ jetzt noch frei, weil
das Teil noch nicht gelegt ist. Das widerspricht der Wahl von $c$ als
\emph{erster} freier Zelle. Also ist $m = c$, und es genügt, jedes Teil nur an dieser
einen Stelle auszuprobieren. Das reduziert die Suche stark, weil nicht jede Position
der Kiste geprüft werden muss.

\subsection{Schritt 5: Backtracking}

\begin{algorithm}[H]
\caption{\code{solvePuzzle} (vereinfacht)}
\begin{algorithmic}[1]
\Require Kiste \textit{box}, Abbildung \textit{typen}: Teiltyp $\to$ freie IDs, Zelle $c$ (optional)
\Ensure \code{true}, wenn die Kiste vollständig gefüllt werden kann
\Function{Solve}{\textit{box}, \textit{typen}, $c$}
  \If{\textit{typen} ist leer} \Return \code{true} \Comment{alle Teile gelegt}
  \EndIf
  \If{$c$ existiert nicht} \Return \code{false} \Comment{kein Platz mehr, aber Teile übrig}
  \EndIf
  \For{\textbf{each} Teiltyp $T$ in \textit{typen}}
    \State nimm eine freie ID $i$ von $T$ aus \textit{typen} heraus
    \For{\textbf{each} Orientierung $o$ von $T$}
      \If{\Call{placeBrick}{\textit{box}, $T_i$, $o$, $c$}} \Comment{passt in die Kiste und überlappt nicht}
        \State $c' \gets$ \Call{nextFreeCell}{$c$}
        \If{\Call{Solve}{\textit{box}, \textit{typen}, $c'$}}
          \State \Return \code{true}
        \EndIf
        \State \Call{removeBrick}{\textit{box}, $T_i$, $o$, $c$} \Comment{Backtracking: Zug zurücknehmen}
      \EndIf
    \EndFor
    \State lege ID $i$ wieder in \textit{typen} zurück
  \EndFor
  \State \Return \code{false}
\EndFunction
\end{algorithmic}
\end{algorithm}

\code{placeBrick} berechnet zuerst alle Zellen, die das Teil in der gewählten
Orientierung ab Position $c$ belegen würde (\code{getCellsForBrick}). Liegt eine
dieser Zellen außerhalb der Kiste oder ist sie schon belegt (auch durch den goldenen
Würfel), wird nichts verändert und \code{false} zurückgegeben. Andernfalls werden alle
Zellen belegt.

Wird \code{true} zurückgegeben, enthält \code{box} die vollständige Lösung. Wird
\code{false} zurückgegeben, wurden alle Möglichkeiten erschöpfend geprüft, und das
Rätsel ist nachweislich unlösbar.

\subsection{Komplexität}

Im schlechtesten Fall wächst die Laufzeit exponentiell mit der Zahl der Teile, denn
das Problem ist ein Packproblem. Die Volumenprüfung, die Zusammenfassung gleicher
Teile, das Entfernen doppelter Orientierungen und das Füllen der ersten freien Zelle
schneiden aber große Teile des Suchbaums ab. In der Praxis sind die Laufzeiten
deshalb gering (siehe Tabelle~\ref{tab:ergebnisse}).

% =====================================================================
\section{Klassendiagramm (OOD)}
% =====================================================================

Die Klassen sind auf zwei Pakete verteilt: \code{core} enthält die Fachlogik ohne
Abhängigkeit zu JavaFX, \code{visualizer} enthält die Oberfläche. So lässt sich der
Kern unabhängig von der Oberfläche testen und wiederverwenden.

Notation: \code{+}~öffentlich, \code{-}~privat, \code{\textasciitilde}~paketsichtbar,
\st{unterstrichen}~statisch. Durchgezogene Pfeile sind Assoziationen (Attribute),
gestrichelte Pfeile Abhängigkeiten (\enquote{benutzt}/\enquote{erzeugt}), Rauten
Kompositionen bzw. Aggregationen.

\subsection{Paket \code{core}}

\begin{figure}[H]
\centering
\resizebox{\textwidth}{!}{%
\begin{tikzpicture}[node distance=1.1cm and 1.8cm]
  % ---------- Zeile 1 ----------
  \umlclass[anchor=north, umlclass=9.4cm]{PuzzleInput}{\stereo{record} PuzzleInput}
    {+ box : Box\\+ bricks : ArrayList<Brick>}
    {+ \st{fromText(text : String) : PuzzleInput}\\
     + \st{fromFile(fileName : String) : PuzzleInput}}
  \umlclass[right=of PuzzleInput.north east, anchor=north west, umlclass=9.4cm]{FileManager}{\stereo{final} FileManager}
    {~}
    {+ \st{readFile(fileName : String) : String}\\
     + \st{writeFile(fileName : String, content : String)}}

  % ---------- Zeile 2 ----------
  \umlclass[below=of PuzzleInput, umlclass=9.4cm]{Puzzle}{\stereo{final} Puzzle}
    {+ \st{orientations : List<Map<Character,Character>{}>}}
    {+ \st{solvePuzzle(box : Box, bricks : ArrayList<Brick>) : boolean}\\
     - \st{solvePuzzle(box, distinct, pos : Optional<Position>) : boolean}\\
     - \st{tryBrick(box, distinct, brick, pos : Position) : boolean}\\
     + \st{distinctBricks(bricks : List<Brick>) : Map<Brick,Set<Integer>{}>}\\
     - \st{unsolvableFromStart(box, bricks) : boolean}\\
     + \st{getOrientation(c : Character, brick : Brick) : int}\\
     + \st{parseBoxString(s : String) : Box}\\
     + \st{parseBrickString(id : int, s : String) : Brick}\\
     + \st{parseMultipleBricksString(s : String) : Brick[]}}
  \umlclass[below=of FileManager, umlclass=9.4cm]{Box}{Box}
    {- placedBricks : Stack<Brick>\\
     - grid : Brick[][][]\\
     - size : int\\
     - goldBrickPosition : Position\\
     - goldBrick : Brick}
    {+ Box(size : int)\\
     + placeBrick(b : Brick, o : Map<..>, p : Position) : boolean\\
     + removeBrick(b : Brick, o : Map<..>, p : Position)\\
     + removeBrick(b : Brick)\\
     - getCellsForBrick(b, o, p) : Position[]\\
     + nextCell(p : Position) : Position\\
     + nextFreeCell(p : Position) : Optional<Position>\\
     + positionOutsideBox(p : Position) : boolean\\
     + isEmpty(p : Position) : boolean\\
     - cellFree(p : Position) : boolean\\
     + printBox() : String\\
     + getGrid() : Brick[][][]\\
     + getSize() : int}

  % ---------- Zeile 3 (unterhalb der höheren Klasse Box) ----------
  \coordinate (row3) at ($(Box.south)+(0,-1.4)$);
  \umlclass[anchor=north west, umlclass=9.4cm, at={(Puzzle.west |- row3)}]{Brick}{\stereo{record} Brick}
    {+ id : int\\+ x : int\\+ y : int\\+ z : int\\
     + orientations : List<Map<Character,Character>{}>}
    {\textasciitilde\ Brick(id, x, y, z : int)\\
     + equals(o : Object) : boolean\\
     + hashCode() : int\\
     - sortedDimensions() : List<Integer>\\
     - \st{getOrientations(x, y, z : int) : List<Map<...>{}>}}
  \umlclass[anchor=north west, umlclass=3.6cm, at={(Box.west |- row3)}]{Position}{\stereo{record} Position}
    {+ x : int\\+ y : int\\+ z : int}{~}
  \umlclass[right=1.2cm of Position.north east, anchor=north west, umlclass=3.6cm]{Triple}{\stereo{record} Triple}
    {+ x : int\\+ y : int\\+ z : int}{~}

  % ---------- Beziehungen ----------
  \draw[dep] (PuzzleInput.east) -- node[lbl, above]{\stereo{use}} (PuzzleInput.east -| FileManager.west);
  \draw[dep] (PuzzleInput) -- node[lbl, right]{\stereo{use}} (Puzzle);
  \draw[assoc] (PuzzleInput.south east) -- node[lbl, pos=0.5]{box 1} (Box.north west);
  \draw[assoc] (PuzzleInput.west) -- ++(-0.7,0) |- node[lbl, pos=0.25]{bricks 0..*} (Brick.west);
  \draw[dep] (Puzzle.east) -- node[lbl, above]{\stereo{use}} (Puzzle.east -| Box.west);
  \draw[dep] (Puzzle.south) -- node[lbl, right]{\stereo{create}} (Puzzle.south |- Brick.north);
  \draw[aggr] (Box.south west) ++(0.6,0) coordinate (bs)
              (bs) -- node[lbl, pos=0.45]{grid 0..*} (Brick.north east);
  \draw[assoc] (Position.north |- Box.south) -- node[lbl, right]{goldBrickPosition 1} (Position.north);
  \draw[dep] (Brick.south) -- ++(0,-0.5) -| node[lbl, pos=0.25, below]{\stereo{use}} (Triple.south);
\end{tikzpicture}}
\caption{Klassendiagramm des Pakets \code{GoldeneMitte.core}}
\label{fig:core}
\end{figure}

\begin{description}
  \item[\code{Puzzle}] Statische Hilfsklasse mit der Fachlogik: Einlesen der
    Textbausteine (\code{parse...}), Volumenprüfung, Gruppierung gleicher Teile und
    der rekursive Backtracking-Algorithmus (\code{solvePuzzle}, \code{tryBrick}).
    Sie enthält außerdem die sechs möglichen Achsen-Permutationen
    (\code{orientations}).
  \item[\code{Box}] Die Kiste als dreidimensionales Gitter. Sie kapselt alle
    Zugriffe auf das Gitter: Teile legen und entfernen, freie Zellen finden, Grenzen
    prüfen und die Kiste als Text ausgeben. Der goldene Würfel wird im Konstruktor
    fest in die Mitte gelegt.
  \item[\code{Brick}] Unveränderlicher Baustein (Java-\code{record}) mit ID, Maßen
    und vorberechneten, eindeutigen Orientierungen. \code{equals}/\code{hashCode}
    vergleichen bewusst nur die sortierten Maße und nicht die ID, damit gleiche Teile
    in der \code{Map} von \code{distinctBricks} zusammenfallen.
  \item[\code{Position}, \code{Triple}] Einfache Werttypen für eine Gitterzelle bzw.
    für die Maße einer Orientierung (zum Entfernen von Duplikaten).
  \item[\code{PuzzleInput}] Bündelt das Ergebnis des Einlesens (Kiste und Teile)
    und bietet mit \code{fromFile}/\code{fromText} \emph{eine} zentrale Stelle zum
    Parsen. Diese wird von \code{Main}, \code{Visualizer} und den Tests gemeinsam
    genutzt.
  \item[\code{FileManager}] Lesen und Schreiben von Textdateien. Ein-/Ausgabefehler
    werden in eine \code{RuntimeException} mit dem absoluten Dateipfad umgewandelt.
\end{description}

\subsection{Paket \code{visualizer} und Einstiegspunkte}

Die Oberfläche folgt dem \textbf{Model-View-Controller}-Muster:
\begin{itemize}
  \item Das \emph{Model} (\code{VoxelModel}) kennt die belegten Zellen je Baustein
        und den aktuell ausgewählten Baustein.
  \item Die \emph{Views} (\code{MainView}, \code{SidebarView}, \code{SceneView})
        stellen dar und melden Benutzereingaben weiter, ändern das Model aber nicht
        selbst.
  \item Die \emph{Controller} (\code{VisualizerController}, \code{CameraController})
        verbinden beide Seiten.
\end{itemize}

\begin{figure}[tbp]
\centering
\resizebox{\textwidth}{!}{%
\begin{tikzpicture}[node distance=1.0cm and 1.2cm]
  % ---------- Spalte 1: Einstieg + Model ----------
  \umlclass[anchor=north, umlclass=7.3cm]{Launcher}{Launcher}
    {~}{+ \st{main(args : String[])}}
  \umlclass[below=of Launcher, umlclass=7.3cm]{VoxelModel}{VoxelModel}
    {+ \st{GOLD\_BRICK\_ID : int = -1}\\
     - gridSize : int\\
     - cellsByBrickId : Map<Integer,List<Position>{}>\\
     - selectedBrick : ObjectProperty<Integer>}
    {+ VoxelModel(box : Box)\\
     + getGridSize() : int\\
     + getBrickIds() : List<Integer>\\
     + getCells(id : int) : List<Position>\\
     + selectedProperty() : ObjectProperty<Integer>\\
     + getSelected() : Integer\\
     + toggleSelection(id : int)\\
     + clearSelection()}
  \umlclass[below=of VoxelModel, umlclass=7.3cm]{Theme}{\stereo{final} Theme}
    {\textasciitilde\ \st{GOLD\_COLOR, OUTLINE\_EDGE, ...: Color}\\
     \textasciitilde\ \st{MAIN\_STYLE, ROW\_STYLE, ...: String}}{~}

  % ---------- Spalte 2: Controller ----------
  \umlclass[right=of Launcher.north east, anchor=north west, umlclass=7.3cm]{Visualizer}{Visualizer {\normalfont\itshape extends Application}}
    {~}{+ start(stage : Stage)}
  \umlclass[below=of Visualizer, umlclass=7.3cm]{VisualizerController}{VisualizerController}
    {- model : VoxelModel\\- view : MainView}
    {+ VisualizerController(model, view)\\
     - connectModelToView()\\
     - connectViewToModel()\\
     - onSceneClicked(brickId : Integer)}
  \umlclass[below=of VisualizerController, umlclass=7.3cm]{CameraController}{CameraController}
    {- \st{ROTATE\_SPEED, ZOOM\_SPEED, DRAG\_THRESHOLD}\\
     - sceneView : SceneView\\
     - onSceneClicked : Consumer<Integer>\\
     - lastMouseX/Y, pressStartX/Y : double\\
     - isDragging : boolean}
    {\textasciitilde\ CameraController(view, onSceneClicked)\\
     - setMouseHandlers()}
  \umlclass[below=of CameraController, umlclass=7.3cm]{BrickStyles}{BrickStyles}
    {- stylesByBrickId : Map<Integer,BrickStyle>}
    {+ getBrickStyle(id : int) : BrickStyle\\
     \textasciitilde\ \st{getColor(id : int) : Color}}
  \umlclass[below=of BrickStyles, umlclass=7.3cm]{BrickStyle}{\stereo{record} BrickStyle}
    {+ color : Color\\+ material : PhongMaterial}{~}

  % ---------- Spalte 3: Views ----------
  \umlclass[right=of Visualizer.north east, anchor=north west, umlclass=7.3cm]{MainView}{MainView {\normalfont\itshape extends BorderPane}}
    {- sidebar : SidebarView\\- sceneView : SceneView}
    {+ MainView(model : VoxelModel)\\
     + getSceneView() : SceneView\\
     + showSelection(selected : Integer)\\
     + setOnBrickClicked(h : IntConsumer)\\
     + setOnResetCameraRequested(h : Runnable)}
  \umlclass[below=of MainView, umlclass=7.3cm]{SidebarView}{SidebarView {\normalfont\itshape extends VBox}}
    {- rows : Map<Integer,HBox>\\- onBrickClicked : IntConsumer}
    {+ showSelection(selected : Integer)\\
     + setOnBrickClicked(h : IntConsumer)\\
     - createRow(id, color, cells) : HBox}
  \umlclass[below=of SidebarView, umlclass=7.3cm]{SceneView}{SceneView {\normalfont\itshape extends StackPane}}
    {- world : Group\\
     - rotX, rotY : Rotate\\
     - camera : PerspectiveCamera\\
     - subScene : SubScene\\
     - voxelNodes : List<VoxelNode>}
    {+ showSelection(selected : Integer)\\
     + findBrickId(n : Node) : Integer\\
     + rotate(dy, dx : double)\\
     + zoom(delta : double)\\
     + resetCamera()\\
     + getSubScene() : SubScene\\
     + setReset(h : Runnable)\\
     - buildFloorGrid(size : int)\\
     - buildVoxels(model : VoxelModel)}
  \umlclass[below=of SceneView, umlclass=7.3cm]{VoxelNode}{VoxelNode}
    {+ brickId : int\\+ cube : Box\\+ outlineBars : List<Box>\\+ group : Group}
    {\textasciitilde\ setOutlineMaterial(m : PhongMaterial)}

  % ---------- Beziehungen ----------
  \draw[dep] (Launcher.east) -- node[lbl, above]{\stereo{launch}} (Launcher.east -| Visualizer.west);
  \draw[dep] (Visualizer.south west) -- node[lbl, pos=0.45]{\stereo{create}} (VoxelModel.north east);
  \draw[dep] (Visualizer) -- node[lbl, right]{\stereo{create}} (VisualizerController);
  \draw[dep] (Visualizer.east) -- node[lbl, above]{\stereo{create}} (Visualizer.east -| MainView.west);
  \draw[assoc] (VisualizerController.west) -- node[lbl, above]{model 1} (VisualizerController.west -| VoxelModel.east);
  \draw[assoc] (VisualizerController.north east) -- node[lbl, pos=0.5]{view 1} (MainView.south west);
  \draw[dep] (VisualizerController) -- node[lbl, right]{\stereo{create}} (CameraController);
  \draw[assoc] (CameraController.south east) -- node[lbl, pos=0.4, below left]{sceneView 1} (SceneView.west);
  \draw[comp] (MainView) -- node[lbl, right]{1} (SidebarView);
  \draw[comp] (MainView.east) -- ++(0.5,0) |- node[lbl, pos=0.75, above]{1} (SceneView.east);
  \draw[comp] (SceneView) -- node[lbl, right]{0..*} (VoxelNode);
  \draw[dep] (SceneView.west |- BrickStyles.east) -- (BrickStyles.east);
  \draw[dep] (SidebarView.west) -- ++(-0.6,0) |- node[lbl, pos=0.3, left]{\stereo{use}} (BrickStyles.east);
  \draw[aggr] (BrickStyles) -- node[lbl, right]{0..*} (BrickStyle);
\end{tikzpicture}}
\caption{Klassendiagramm des Pakets \code{GoldeneMitte.visualizer} mit den Einstiegspunkten}
\label{fig:visualizer}
\end{figure}

\begin{description}
  \item[\code{Launcher}, \code{Visualizer}] \code{Launcher.main} startet die
    JavaFX-Anwendung \code{Visualizer}. Diese liest das Rätsel ein, löst es, speichert
    die Lösung und baut dann Model, View und Controller auf.
  \item[\code{Main}] (nicht abgebildet) ist die Konsolenvariante ohne Oberfläche.
    Sie benutzt nur \code{PuzzleInput}, \code{Puzzle} und \code{Box}.
  \item[\code{VoxelModel}] Übersetzt das \code{Brick[][][]}-Gitter der \code{Box}
    einmalig in eine Abbildung \emph{Baustein-ID $\to$ Liste belegter Zellen}. Die
    Auswahl wird als JavaFX-\code{ObjectProperty} gehalten, damit sich Views über
    Änderungen benachrichtigen lassen können (Beobachter-Muster).
  \item[\code{VisualizerController}] Registriert einen Listener auf die Auswahl im
    Model und leitet Änderungen an \code{MainView.showSelection} weiter. Umgekehrt
    gibt er Klicks aus der View als \code{toggleSelection}/\code{clearSelection} an
    das Model weiter.
  \item[\code{CameraController}] Wertet die Mausereignisse der 3D-Szene aus. Ziehen
    dreht die Kamera, das Mausrad zoomt. Ein einfacher Klick ohne Ziehen (Schwelle
    \code{DRAG\_THRESHOLD}) wählt den getroffenen Baustein aus.
  \item[\code{MainView}] Container (\code{BorderPane}) mit der Seitenleiste links und
    der 3D-Szene in der Mitte. Er bietet dem Controller eine schmale Schnittstelle an.
  \item[\code{SidebarView}] Liste aller Bausteine mit Farbe, ID und Zellanzahl.
  \item[\code{SceneView}] Baut die 3D-Szene aus einem \code{VoxelNode} je Zelle auf,
    zeichnet ein Bodengitter und verwaltet Kamera, Drehung und Zoom.
  \item[\code{VoxelNode}] Eine Zelle in der Szene: ein Würfel plus zwölf
    Kantenbalken als Umriss.
  \item[\code{BrickStyles}, \code{BrickStyle}, \code{Theme}] Farbgebung: Jeder ID
    wird über einen Farbton-Versatz von $137^\circ$ eine gut unterscheidbare Farbe
    zugeordnet. Der goldene Würfel ist immer goldfarben. \code{Theme} sammelt alle
    Farb- und Stilkonstanten an einer Stelle.
\end{description}

% =====================================================================
\section{Tests}
% =====================================================================

\subsection{Vorgehen}

Die Fachlogik im Paket \code{core} wird mit \textbf{JUnit~5} automatisiert getestet.
Die Tests liegen unter \file{src/test/java/GoldeneMitte/core/} und werden mit
\begin{lstlisting}
mvn test
\end{lstlisting}
oder direkt in IntelliJ ausgeführt. Beim Schreiben der Tests wurden folgende
Grundsätze beachtet:
\begin{itemize}
  \item \textbf{Arrange -- Act -- Assert:} Jeder Test baut zuerst die Eingabe auf,
        führt dann genau eine Aktion aus und prüft danach das Ergebnis.
  \item \textbf{Ein Szenario pro Test} und sprechende Testnamen
        (z.\,B. \code{placeBrick\_\allowbreak overlappingBrick\_\allowbreak returnsFalse}). So zeigt schon
        der Name eines fehlgeschlagenen Tests, was kaputt ist.
  \item \textbf{Parametrisierte Tests} (\code{@ParameterizedTest},
        \code{@ValueSource}) prüfen viele ungültige Eingaben mit einem einzigen
        Testfall.
  \item \textbf{Grenzfälle:} Kistengröße 1, gerade und negative Größen, Zellen am
        Rand, Teile, die aus der Kiste ragen, Volumen um genau eins zu groß bzw. zu
        klein.
  \item \textbf{Isolation:} Dateitests schreiben mit \code{@TempDir} in ein
        temporäres Verzeichnis, damit die echten Projektdateien unberührt bleiben.
        Jeder Test bekommt über \code{@BeforeEach} eine frische Kiste.
  \item \textbf{Den echten Code testen:} Die Tests der Beispieldateien laden die
        Dateien über dieselbe Methode \code{PuzzleInput.fromFile}, die auch die
        Anwendung benutzt, statt die Parse-Logik nachzubauen.
\end{itemize}

\subsection{Übersicht der Testklassen}

\begin{center}
\begin{tabularx}{\textwidth}{@{}lcX@{}}
\toprule
\textbf{Testklasse} & \textbf{Tests} & \textbf{Geprüft wird} \\
\midrule
\code{BoxTest} & 20 &
  Konstruktor (gültige/ungültige Größen), goldener Würfel in der Mitte,
  Bereichsprüfung, \code{nextCell}/\code{nextFreeCell} inkl. Zeilen-/Ebenenumbruch
  und Überspringen der Mitte, Legen (passt, überlappt, ragt heraus, auf den
  goldenen Würfel) und Entfernen von Teilen, Textausgabe. \\
\code{FileManagerTest} & 3 &
  Schreiben und Wiederlesen, fehlende Datei, fehlendes Verzeichnis. \\
\code{PuzzleTest} & 19 &
  Parsen von Kisten- und Bausteinzeilen (gültig/ungültig, Windows-Zeilenenden),
  Lösen des Beispiels aus der Aufgabenstellung, Volumen zu groß/zu klein,
  richtiges Volumen mit nicht passendem Teil, Wiederholbarkeit. \\
\code{InvalidPuzzleFilesTest} & 2 &
  Die Dateien \file{raetsel\_nicht\_lösbar.txt} und
  \file{raetsel\_unvollständig.txt} werden vollständig eingelesen, als unlösbar
  erkannt, und die Kiste bleibt bis auf den goldenen Würfel leer. \\
\midrule
\textbf{Summe} & \textbf{44} & \\
\bottomrule
\end{tabularx}
\end{center}

\subsection{Testergebnis}

Beim letzten Testlauf (JDK~25.0.1) liefen \textbf{44 Tests}. Davon waren \textbf{41
erfolgreich} und \textbf{3 fehlgeschlagen}. Alle drei Fehlschläge gehören zum
parametrisierten Test \code{PuzzleTest.\allowbreak parseBrickString\_\allowbreak invalid\_\allowbreak throws}. Sie
zeigen echte Lücken in der Eingabeprüfung von \code{Puzzle.parseBrickString}:

\begin{center}
\begin{tabularx}{\textwidth}{@{}lX@{}}
\toprule
\textbf{Eingabe} & \textbf{Beobachtetes Verhalten} \\
\midrule
\code{"0 1 1"} & Wird akzeptiert. Die Prüfung verlangt nur
  \code{>= 0}, Kantenlängen von~0 werden also nicht abgelehnt. \\
\code{"1 2"} & Wirft \code{ArrayIndexOutOfBoundsException} statt
  \code{IllegalArgumentException}, weil die Anzahl der Werte nicht geprüft wird. \\
\code{"1 2 3 4"} & Wird akzeptiert, der überzählige Wert wird stillschweigend
  ignoriert. \\
\bottomrule
\end{tabularx}
\end{center}

Für die mitgelieferten Rätseldateien spielen diese Fälle keine Rolle. Sie sollten
aber behoben werden, indem \code{parseBrickString} genau drei Werte $\geq 1$
verlangt.

\subsection{Anwendung auf alle Beispiele}
\label{sec:ergebnisse}

Das Programm wurde auf alle Beispieldateien angewendet. Die Laufzeiten stammen aus
einem einzelnen Lauf auf einem Entwicklungsrechner (JDK~25, alle Rätsel nacheinander
in derselben JVM) und sind nur als Größenordnung zu verstehen.

\begin{table}[H]
\centering
\begin{tabular}{@{}llrrll@{}}
\toprule
\textbf{Datei} & \textbf{Kiste} & \textbf{Teile} & \textbf{Volumen} & \textbf{Ergebnis} & \textbf{Laufzeit} \\
\midrule
\file{raetsel1.txt} & $3^3$ & 8  & 26 / 26   & lösbar       & ca.\ 20\,ms \\
\file{raetsel2.txt} & $3^3$ & 6  & 26 / 26   & lösbar       & < 10\,ms \\
\file{raetsel3.txt} & $3^3$ & 6  & 26 / 26   & nicht lösbar & < 10\,ms \\
\file{raetsel4.txt} & $5^3$ & 20 & 124 / 124 & lösbar       & ca.\ 3\,s \\
\file{raetsel5.txt} (Dirk) & $5^3$ & 16 & 124 / 124 & lösbar & ca.\ 0{,}5\,s \\
\file{raetsel\_nicht\_lösbar.txt}   & $3^3$ & 8 & 28 / 26 & nicht lösbar & < 1\,ms \\
\file{raetsel\_unvollständig.txt}   & $3^3$ & 7 & 22 / 26 & nicht lösbar & < 1\,ms \\
\bottomrule
\end{tabular}
\caption{Ergebnisse für alle Beispiele (Volumen: Summe der Teile / freie Zellen $n^3-1$)}
\label{tab:ergebnisse}
\end{table}

\paragraph{Warum ist \file{raetsel3.txt} unlösbar?}
Das Volumen stimmt, deshalb wird hier die vollständige Suche durchlaufen. Sie
scheitert aus einem einfachen geometrischen Grund: Ein $2\times2\times2$-Würfel
belegt in einer $3\times3\times3$-Kiste in jeder Achse die Koordinaten $\{0,1\}$ oder
$\{1,2\}$. Er enthält also immer die Koordinate~1 in allen drei Achsen und damit die
Mittelzelle $(1,1,1)$, in der der goldene Würfel liegt.

\paragraph{Warum sind die beiden zusätzlichen Dateien unlösbar?}
In \file{raetsel\_nicht\_lösbar.txt} sind die Teile zusammen zu groß (28 statt 26
Zellen). Außerdem passt das Teil $1\times1\times4$ in keiner Lage in eine Kiste der
Kantenlänge~3. In \file{raetsel\_unvollständig.txt} fehlen Teile, sodass vier Zellen
leer bleiben würden. Beide Fälle erkennt die Volumenprüfung sofort.

\subsubsection*{Lösung für Dirks Rätsel (\file{raetsel5.txt})}

Darstellung Ebene für Ebene ($z = 0 \dots 4$). In jeder Ebene stehen die Zeilen für
$y = 0 \dots 4$ und die Spalten für $x = 0 \dots 4$. Die Zahl ist die ID des Teils
(0-basiert), \colorbox{goldcol}{\textbf{G}} ist der goldene Würfel. Die IDs 0 und
13--15 sind die $1\times1\times1$-Würfel, 1--6 die $2\times2\times3$-Quader und 7--12
die $1\times2\times4$-Quader.

\begin{center}\footnotesize\setlength{\tabcolsep}{3pt}
\begin{tabular}[t]{@{}ccccc@{}}
\multicolumn{5}{c}{\textbf{Ebene 0}}\\
0 & 7 & 7 & 7 & 7\\
8 & 7 & 7 & 7 & 7\\
8 & 1 & 1 & 2 & 2\\
8 & 1 & 1 & 2 & 2\\
8 & 1 & 1 & 9 & 9\\
\end{tabular}\hspace{0.9em}
\begin{tabular}[t]{@{}ccccc@{}}
\multicolumn{5}{c}{\textbf{Ebene 1}}\\
10 & 10 & 3 & 3 & 3\\
8 & 13 & 3 & 3 & 3\\
8 & 1 & 1 & 2 & 2\\
8 & 1 & 1 & 2 & 2\\
8 & 1 & 1 & 9 & 9\\
\end{tabular}\hspace{0.9em}
\begin{tabular}[t]{@{}ccccc@{}}
\multicolumn{5}{c}{\textbf{Ebene 2}}\\
10 & 10 & 3 & 3 & 3\\
4 & 4 & 3 & 3 & 3\\
4 & 4 & \gold & 2 & 2\\
5 & 5 & 5 & 2 & 2\\
5 & 5 & 5 & 9 & 9\\
\end{tabular}\hspace{0.9em}
\begin{tabular}[t]{@{}ccccc@{}}
\multicolumn{5}{c}{\textbf{Ebene 3}}\\
10 & 10 & 6 & 6 & 11\\
4 & 4 & 6 & 6 & 11\\
4 & 4 & 6 & 6 & 11\\
5 & 5 & 5 & 14 & 11\\
5 & 5 & 5 & 9 & 9\\
\end{tabular}\hspace{0.9em}
\begin{tabular}[t]{@{}ccccc@{}}
\multicolumn{5}{c}{\textbf{Ebene 4}}\\
10 & 10 & 6 & 6 & 11\\
4 & 4 & 6 & 6 & 11\\
4 & 4 & 6 & 6 & 11\\
12 & 12 & 12 & 12 & 11\\
12 & 12 & 12 & 12 & 15\\
\end{tabular}\hspace{0.9em}
\end{center}

\subsubsection*{Lösung für \file{raetsel4.txt}}

\begin{center}\footnotesize\setlength{\tabcolsep}{3pt}
\begin{tabular}[t]{@{}ccccc@{}}
\multicolumn{5}{c}{\textbf{Ebene 0}}\\
19 & 16 & 5 & 1 & 1\\
19 & 16 & 6 & 1 & 1\\
19 & 17 & 18 & 18 & 18\\
14 & 17 & 7 & 2 & 2\\
14 & 8 & 9 & 2 & 2\\
\end{tabular}\hspace{0.9em}
\begin{tabular}[t]{@{}ccccc@{}}
\multicolumn{5}{c}{\textbf{Ebene 1}}\\
19 & 16 & 5 & 1 & 1\\
19 & 16 & 6 & 1 & 1\\
19 & 17 & 18 & 18 & 18\\
14 & 17 & 7 & 2 & 2\\
14 & 8 & 9 & 2 & 2\\
\end{tabular}\hspace{0.9em}
\begin{tabular}[t]{@{}ccccc@{}}
\multicolumn{5}{c}{\textbf{Ebene 2}}\\
19 & 16 & 5 & 10 & 11\\
19 & 16 & 6 & 10 & 11\\
19 & 17 & \gold & 10 & 11\\
14 & 17 & 7 & 10 & 11\\
14 & 8 & 9 & 10 & 11\\
\end{tabular}\hspace{0.9em}
\begin{tabular}[t]{@{}ccccc@{}}
\multicolumn{5}{c}{\textbf{Ebene 3}}\\
12 & 0 & 5 & 3 & 3\\
12 & 0 & 6 & 3 & 3\\
12 & 0 & 15 & 15 & 15\\
12 & 0 & 7 & 4 & 4\\
12 & 8 & 9 & 4 & 4\\
\end{tabular}\hspace{0.9em}
\begin{tabular}[t]{@{}ccccc@{}}
\multicolumn{5}{c}{\textbf{Ebene 4}}\\
13 & 0 & 5 & 3 & 3\\
13 & 0 & 6 & 3 & 3\\
13 & 0 & 15 & 15 & 15\\
13 & 0 & 7 & 4 & 4\\
13 & 8 & 9 & 4 & 4\\
\end{tabular}\hspace{0.9em}
\end{center}

\subsubsection*{Lösung für \file{raetsel1.txt}}

\begin{center}\footnotesize\setlength{\tabcolsep}{3pt}
\begin{tabular}[t]{@{}ccc@{}}
\multicolumn{3}{c}{\textbf{Ebene 0}}\\
4 & 5 & 5\\
4 & 6 & 6\\
0 & 6 & 6\\
\end{tabular}\hspace{0.9em}
\begin{tabular}[t]{@{}ccc@{}}
\multicolumn{3}{c}{\textbf{Ebene 1}}\\
4 & 5 & 5\\
4 & \gold & 7\\
0 & 2 & 7\\
\end{tabular}\hspace{0.9em}
\begin{tabular}[t]{@{}ccc@{}}
\multicolumn{3}{c}{\textbf{Ebene 2}}\\
1 & 1 & 1\\
3 & 3 & 7\\
0 & 2 & 7\\
\end{tabular}\hspace{0.9em}
\end{center}

\subsubsection*{Lösung für \file{raetsel2.txt}}

\begin{center}\footnotesize\setlength{\tabcolsep}{3pt}
\begin{tabular}[t]{@{}ccc@{}}
\multicolumn{3}{c}{\textbf{Ebene 0}}\\
0 & 0 & 0\\
0 & 0 & 0\\
0 & 0 & 0\\
\end{tabular}\hspace{0.9em}
\begin{tabular}[t]{@{}ccc@{}}
\multicolumn{3}{c}{\textbf{Ebene 1}}\\
4 & 4 & 4\\
2 & \gold & 3\\
5 & 5 & 5\\
\end{tabular}\hspace{0.9em}
\begin{tabular}[t]{@{}ccc@{}}
\multicolumn{3}{c}{\textbf{Ebene 2}}\\
1 & 1 & 1\\
1 & 1 & 1\\
1 & 1 & 1\\
\end{tabular}\hspace{0.9em}
\end{center}

% =====================================================================
\section{Fazit}
% =====================================================================

Das Programm erfüllt die Aufgabenstellung: Es füllt Dirks Kiste vollständig mit dem
goldenen Würfel in der Mitte und findet die Lösung in etwa einer halben Sekunde. Es
funktioniert auch für andere Kisten mit ungerader Kantenlänge und beliebige Mengen von
Quadern. Für unlösbare Rätsel meldet es zuverlässig, dass es keine Lösung gibt.
Dieses Ergebnis ist korrekt, weil das Backtracking alle Möglichkeiten erschöpfend
prüft.

\paragraph{Was gut funktioniert hat.}
\begin{itemize}
  \item Die Kombination aus Volumenprüfung, Gruppierung gleicher Teile,
        doppelfreien Orientierungen und dem Füllen der jeweils ersten freien Zelle
        hält den Suchraum klein. Ohne diese Maßnahmen wäre schon das
        $5\times5\times5$-Rätsel praktisch nicht lösbar.
  \item Die strikte Trennung von Fachlogik (\code{core}) und Oberfläche
        (\code{visualizer}, MVC) hat die automatisierten Tests erst möglich gemacht
        und hält die Klassen klein und verständlich.
  \item Die 3D-Ansicht macht Lösungen deutlich leichter nachvollziehbar als die
        Textausgabe, vor allem durch die Auswahl einzelner Bausteine.
\end{itemize}

\paragraph{Grenzen und Ausblick.}
\begin{itemize}
  \item Die Eingabedatei ist fest im Quelltext eingetragen. Sinnvoll wäre die Auswahl
        per Kommandozeilenparameter oder Dateidialog in der Oberfläche.
  \item Die Eingabeprüfung ist unvollständig. Die Anzahl $k$ in der zweiten Zeile
        wird nicht mit der tatsächlichen Zahl der Teile verglichen, und
        \code{parseBrickString} akzeptiert Kantenlängen von~0 sowie zu viele Werte
        (siehe fehlgeschlagene Tests). Auch nicht würfelförmige Kisten werden nicht
        gemeldet, sondern es wird nur die erste Kantenlänge verwendet.
  \item \code{Box.removeBrick(brick, orientation, position)} gibt die Zellen frei,
        entfernt den Baustein aber nicht vom Stapel \code{placedBricks}. Nach dem
        Backtracking kann dieser Stapel deshalb Teile enthalten, die nicht mehr in
        der Kiste liegen. Für das Ergebnis spielt das derzeit keine Rolle, weil der
        Stapel nicht ausgewertet wird.
  \item Für größere Kisten wächst die Laufzeit exponentiell. Zusätzliche
        Abbruchkriterien würden helfen, etwa das Erkennen abgeschlossener
        Hohlräume, deren Volumen sich mit den übrigen Teilen nicht füllen lässt. Ein
        Ansatz dafür ist im Code bereits auskommentiert vorbereitet.
  \item Die Textausgabe könnte an das in der Aufgabenstellung vorgeschlagene Format
        (Ebenen nebeneinander, IDs ab~1) angepasst werden.
\end{itemize}

\end{document}
